\documentclass[10pt,a4paper]{article}
\usepackage[margin=18mm,headheight=14pt,headsep=7mm,footskip=10mm]{geometry}
\usepackage[T1]{fontenc}
\usepackage{lmodern,microtype,amsmath,amssymb,booktabs,tabularx,array,xcolor,tikz,fancyhdr,enumitem,hyperref}
\usetikzlibrary{arrows.meta,positioning,calc}
\definecolor{navy}{HTML}{17334A}
\definecolor{teal}{HTML}{197C80}
\definecolor{pale}{HTML}{EDF6F6}
\definecolor{slate}{HTML}{536474}
\definecolor{linegray}{HTML}{D5DFE5}
\definecolor{amber}{HTML}{996519}
\hypersetup{colorlinks=true,urlcolor=teal,linkcolor=navy,pdftitle={FININ: prototype reproduction progress report},pdfauthor={Research progress report}}
\pagestyle{fancy}\fancyhf{}
\fancyhead[L]{\small\sffamily\color{slate}FININ / REPRODUCTION STUDY}
\fancyhead[R]{\small\sffamily\color{slate}4 OCTOBER 2026}
\fancyfoot[L]{\footnotesize\sffamily Research progress report}
\fancyfoot[R]{\footnotesize\sffamily\thepage\ / 5}
\renewcommand{\headrulewidth}{0.3pt}
\setlength{\parindent}{0pt}\setlength{\parskip}{6pt}
\setlength{\tabcolsep}{5pt}\renewcommand{\arraystretch}{1.25}
\setlist[itemize]{leftmargin=*,itemsep=3pt,topsep=3pt}
\newcolumntype{Y}{>{\raggedright\arraybackslash}X}
\newcommand{\topic}[2]{\vspace{3pt}{\small\sffamily\bfseries\color{teal}#1}\par\vspace{1pt}{\LARGE\bfseries\color{navy}#2}\par\vspace{3pt}}
\newcommand{\sub}[1]{\par\vspace{5pt}{\large\bfseries\color{navy}#1}\par}
\newcommand{\note}[1]{\par\noindent\colorbox{pale}{\parbox{\dimexpr\linewidth-2\fboxsep}{\vspace{4pt}#1\vspace{4pt}}}\par}
\newcommand{\code}[1]{{\small\ttfamily\detokenize{#1}}}
\tikzset{box/.style={draw=linegray,rounded corners=2pt,fill=white,align=center,font=\scriptsize,smooth,text=navy,inner sep=4pt},core/.style={box,fill=pale,draw=teal},arr/.style={-{Stealth[length=1.8mm]},draw=slate,line width=.6pt}}

% One conceptual architecture, instantiated with the paper and prototype inputs.
\newcommand{\architecture}[5]{
\begin{tikzpicture}[x=1cm,y=1cm]
\node[font=\small\sffamily\bfseries,text=navy] at (1.8,.55) {#1};
\node[box,text width=2.65cm,minimum height=1.05cm] (market) at (0,-.4) {#2};
\node[box,text width=2.65cm,minimum height=1.05cm] (news) at (3.6,-.4) {#3};
\node[core,text width=2.65cm,minimum height=.95cm] (mf) at (0,-1.95) {Text projection\\+ numeric MLP\\Market fusion MLP};
\node[core,text width=2.65cm,minimum height=.95cm] (nf) at (3.6,-1.95) {Text projection\\+ numeric MLP\\News fusion MLP};
\draw[arr] (market)--(mf);\draw[arr] (news)--(nf);
\node[core,text width=2.65cm,minimum height=.7cm] (sa) at (3.6,-3.25) {News self-attention\\headline interactions};
\draw[arr] (nf)--(sa);
\node[core,text width=2.65cm,minimum height=.9cm] (ca) at (3.6,-4.7) {Market-query attention\\softmax news weights};
\draw[arr] (sa)--node[right,font=\tiny,text=slate]{keys}(ca);
\draw[arr] (mf)--(0,-4.7)--node[above,font=\tiny,text=slate]{market query}(ca);
\node[core,text width=2.65cm,minimum height=.75cm] (sum) at (3.6,-6.1) {Weighted sum of\\refined news vectors};
\draw[arr] (ca)--(sum);
\draw[arr] (sa.east)--++(.24,0)|-(sum.east);
\node[core,text width=5.5cm,minimum height=.85cm] (pred) at (1.8,-7.6) {#4};
\draw[arr] (sum)--(3.6,-6.85)--(pred.north -| sum);
\draw[arr] (0,-4.7)--(pred.north -| mf);
\node[box,text width=5.5cm,minimum height=.55cm] (out) at (1.8,-8.9) {#5};
\draw[arr] (pred)--(out);
\end{tikzpicture}}

\begin{document}
\topic{01 / RESEARCH OBJECTIVE}{FININ: a working reduced prototype}
{\large Reproducing news interactions for next-session market prediction}

\textbf{Paper.} Wang, Cohen and Ma (2024), \emph{Modeling News Interactions and Influence for Financial Market Prediction}, Findings of EMNLP, pp.\ 3302--3314 [1].

\textbf{Research question.} Can a model improve a market forecast by first relating headlines to one another, then weighting them using the current market state? FININ combines each headline's meaning and sentiment, rather than reducing the entire news day to one sentiment score.

\note{\textbf{Progress assessment: sufficient for a working architectural prototype.} The data-to-prediction pipeline runs, all 18 trained checkpoints replay correctly, and the core fusion and attention mechanisms are implemented. The published performance gains and financial conclusions have \textbf{not} been reproduced.}

\sub{Paper architecture versus our implementation}
\begin{center}
\begin{minipage}[t]{.48\linewidth}\centering
\architecture{PAPER / FININ}{Index description +\\constituent names;\\six market fields}{TRNA headlines +\\instrument-related\\sentiment probabilities}{Market + news representations over $t$ days\\MLP predictor; $t\in\{1,3,5,10,20\}$}{Next-session index direction}
\end{minipage}\hfill
\begin{minipage}[t]{.48\linewidth}\centering
\architecture{OURS / REDUCED FININ}{Fixed SPY description;\\six transformed\\market channels}{Up to 16 delayed\\NIFTY headlines +\\TinyBERT probabilities}{One market vector + one news summary\\$128\to64\to1$ MLP; sigmoid probability}{Next-session SPY direction; threshold 0.5}
\end{minipage}
\end{center}
{\footnotesize\color{slate}\textbf{Figure 1.} Conceptual redraw of paper Figure 2 and the implemented data flow. Teal blocks retain the paper's functional structure. Market features also reach the predictor directly; refined news supplies both attention keys and the weighted-sum values.}

\textbf{What is shared.} Within each model, market text and headline text use the same frozen language model and learned projection. Market and news numeric/fusion MLPs remain separate. The paper selects RoBERTa for its main comparisons; ours freezes BGE-small and trains 91,905 downstream parameters.

\newpage
\topic{02 / DATA AND FORECAST TIMING}{Same task family, different evidence}
\begin{tabularx}{\linewidth}{@{}p{2.5cm}YY@{}}
\toprule
\textbf{Dimension} & \textbf{Paper [1]} & \textbf{Our prototype [2; local audit]}\\\midrule
Market & S\&P 500 and NASDAQ 100 indices & SPY ETF, a proxy for the S\&P 500; one market\\
Period & 2003--2018 & News dates: Jan.\ 2010--Sept.\ 2020; last target: 23 Sept.\ 2020\\
News & TRNA: over 2.7 million items; 24--1,953 per day & NIFTY: 436,570 raw headline occurrences over 2,111 dated rows; 33,673 selected slots, 32,286 unique texts\\
Selection & Aims to use all available daily news & Within-day exact deduplication; deterministic SHA-256 ordering; cap of 16\\
Sentiment & Positive/neutral/negative scores about mentioned financial instruments & Three text-sentiment probabilities from a frozen finance TinyBERT; a proxy, not Reuters analytics\\
Market inputs & Open, high, low, close, adjusted close, volume; description and constituent names & Five relative-price channels + $\log(1+\text{volume})$; training-only standardization; fixed description, no constituent list\\
Text encoding & Compares RoBERTa, FinBERT, BGE, Llama2; main results use RoBERTa & BGE-small-en-v1.5: 384 dimensions, first-token pooling, no vector normalization; 128-token cap\\
Evaluation & Ten 500-day windows, stepped by 391 days; chronological 80/10/10 split per window & One chronological source split; four boundary-crossing outcomes removed; three training seeds\\
Scope & 1/3/5/10/20 input days; accuracy, PnL, Sharpe ratio & One input day; accuracy, balanced accuracy, log loss; financial metrics deferred\\
\bottomrule
\end{tabularx}\par\vspace{3pt}
{\footnotesize Counts of raw occurrences, selected slots and unique texts measure different things. NIFTY here denotes the \emph{News-Informed Financial Trend Yield} dataset, not India's NIFTY index.}

\sub{Forecast time is explicit}
Let $n$ be the news date, $d$ its next trading session, and $q$ the session after $d$:
\[
\underbrace{R_n}_{\text{previous-session news}} +
\underbrace{M_d}_{\text{prices through forecast close}}
\quad\longrightarrow\quad
\widehat{P}(C_q>C_d),\qquad y_d=\mathbf{1}[C_q>C_d].
\]
\begin{center}
\begin{tikzpicture}
\node[box,text width=4cm,minimum height=1cm] (a) at (0,0) {\textbf{6 January 2010}\\News bucket selected};
\node[core,text width=4cm,minimum height=1cm] (b) at (5.5,0) {\textbf{7 January: forecast close}\\SPY close = \$114.19};
\node[box,text width=4cm,minimum height=1cm] (c) at (11,0) {\textbf{8 January: outcome}\\Close = \$114.57; target = up};
\draw[arr] (a)--(b);\draw[arr] (b)--(c);
\end{tikzpicture}
\end{center}
We delay the news because individual publication times are unavailable. This reduces timing risk under correct date buckets; it cannot prove historical availability. The horizon remains \textbf{one trading session}, distinct from the input lookback. Source prompts, answers and future prices never enter the feature tensors.

\begin{tabularx}{\linewidth}{@{}Yrrrr@{}}
\toprule
\textbf{Retained split} & \textbf{Dates} & \textbf{News slots} & \textbf{Up targets} & \textbf{Up share}\\\midrule
Training & 1,475 & 23,592 & 807 & 54.7\%\\
Validation & 315 & 5,040 & 172 & 54.6\%\\
Test & 317 & 5,041 & 184 & 58.0\%\\\bottomrule
\end{tabularx}\par\vspace{3pt}
\textbf{Coverage shift.} Raw median headlines/day fall from 78 to 54 to 26 across the splits. The cap controls computation, but cannot equalize relevance or restore omitted stories. All 2,111 supplied NIFTY returns match the next-session SPY return; targets are nevertheless rebuilt for our delayed-news schedule.

\newpage
\topic{03 / EXPERIMENTAL EVIDENCE}{A completed experiment with a negative result}
\textbf{Fixed protocol.} Six trained variants, seeds 17/42/73; identical dates, selected headlines and frozen features. Width 64; two affine layers per MLP; one attention head per layer; predictor dropout 0.1. Adam, learning rate 0.001, batch size 16, binary cross-entropy, gradient clipping at 1.0, up to 30 epochs; early stopping after five epochs without sufficient validation-loss improvement. Checkpoints use validation loss; direction threshold is 0.5.

\sub{Held-out results: 317 dates}
\begin{tabularx}{\linewidth}{@{}Yrrr@{}}
\toprule
\textbf{Model / reference} & \textbf{Accuracy \%} & \textbf{Balanced acc.\ \%} & \textbf{Log loss}\\\midrule
Always up (direction only) & 58.04 & 50.00 & ---\\
Training-prior probability$^{\dagger}$ & 58.04 & 50.00 & 0.6824\\\midrule
Prices only & 58.25 $\pm$ 1.31 & 50.91 $\pm$ 1.50 & 0.6885 $\pm$ 0.0053\\
Prices + mean sentiment & 56.89 $\pm$ 0.48 & 49.91 $\pm$ 0.64 & 0.6884 $\pm$ 0.0053\\
Mean-pooled fused news & 58.04 $\pm$ 0.00 & 50.00 $\pm$ 0.00 & 0.6823 $\pm$ 0.0007\\
No news self-attention & 58.04 $\pm$ 0.00 & 50.00 $\pm$ 0.00 & 0.6816 $\pm$ 0.0008\\
No sentiment input & 58.04 $\pm$ 0.00 & 50.00 $\pm$ 0.00 & 0.6814 $\pm$ 0.0004\\
\textbf{Reduced FININ} & \textbf{58.36 $\pm$ 0.55} & \textbf{50.41 $\pm$ 0.71} & \textbf{0.6816 $\pm$ 0.0011}\\
\bottomrule
\end{tabularx}\par\vspace{3pt}
{\footnotesize Mean $\pm$ sample standard deviation across three seeds; not a confidence interval over future periods. Values are regenerated from saved predictions in the report build. $^{\dagger}$Added during this audit: always output $p=807/1475=0.5471$, fitted only on training labels. Original trained runs are unchanged.}

\note{\textbf{Interpretation.} Reduced FININ is only 0.32 percentage points above always-up accuracy. Its balanced accuracy is approximately chance. A constant training-prior probability also nearly matches its log loss. This experiment provides no convincing evidence of an attention benefit on the substituted dataset.}

\sub{Why accuracy alone is misleading}
\begin{tabularx}{\linewidth}{@{}Yrrrr@{}}
\toprule
\textbf{Full-model seed} & \textbf{Predict up} & \textbf{Predict down} & \textbf{Correct down} & \textbf{Balanced acc.}\\\midrule
17 & 317 & 0 & 0 / 133 & 50.00\%\\
42 & 317 & 0 & 0 / 133 & 50.00\%\\
73 & 312 & 5 & 4 / 133 & 51.23\%\\\bottomrule
\end{tabularx}\par\vspace{3pt}
There are 184 up days and 133 non-up days. Always predicting up therefore earns 58.04\% accuracy while recovering none of the non-up days. Balanced accuracy averages recall across the two classes. Log loss assesses probabilities; assigning probability 1 to every day creates an extreme penalty on non-up days, so the train-prior reference is the useful probability comparison.

\sub{What each comparison tests}
\textbf{Mean pooling} removes both attention mechanisms but retains the fusion pipeline. \textbf{No self-attention} retains market-query weighting. \textbf{No sentiment} replaces the three sentiment inputs with zeros while preserving the network shape. The two numeric baselines are simpler independent models; they are not exact replicas of the paper's ablations.

\textbf{Connection to the publication.} The paper reports its strongest Sharpe ratios at 20 input days: 1.772 on S\&P 500 and 1.449 on NASDAQ 100 [1, Table 2]. These are contextual reference values, not benchmarks for the numbers above: market, data, timing, lookback and evaluation all differ. Our classification result does not confirm or refute that financial-performance claim.

\newpage
\topic{04 / IMPLEMENTATION AND VERIFICATION}{The code skeleton and evidence behind it}
\textbf{End-to-end workflow.} Acquire pinned sources $\to$ audit dates and returns $\to$ build aligned examples $\to$ cache frozen features $\to$ train baselines/ablations $\to$ export predictions and presentation material.

\begin{tabularx}{\linewidth}{@{}p{6.1cm}Y@{}}
\toprule
\textbf{Code area} & \textbf{Role and connection to FININ}\\\midrule
\code{src/data/}\newline\code{nifty.py, prices.py, examples.py} & Load verified snapshots; select headlines; align the trading calendar; derive labels; purge split boundaries and standardize inputs.\\
\code{src/features/cache.py}\newline\code{src/training_data.py} & Frozen BGE and sentiment extraction; cache provenance and hashes; join arrays to dates; pad to 16 headlines and return masks.\\
\code{src/models/encoders.py} & Shared $384\to64$ text projection; distinct $6\to64$ market and $3\to64$ sentiment MLPs; separate $128\to64$ fusion MLPs. Paper Eqs.\ 1--5.\\
\code{src/models/attention.py} & One news self-attention layer, then a market query over refined news keys; masked softmax and weighted sum. Paper Section 4.2 / Eq.\ 6.\\
\code{src/models/finin.py}\newline\code{src/models/baselines.py} & Concatenate market/news summaries; predict one logit; expose weights and structural variants. Reduced Eq.\ 7 with $t=1$.\\
\code{src/training.py}\newline\code{src/evaluation.py} & Seeded optimization, validation selection, checkpoint replay; probabilities, class recalls, confusion matrix and log loss.\\
\code{scripts/01_...}\,through\,\code{06_...}\newline\code{scripts/08_verify_prototype.py} & Ordered execution and offline export; new non-destructive audit reconstructs data and replays every checkpoint.\\
\bottomrule
\end{tabularx}\par\vspace{3pt}

\sub{The central computation}
\[
\widetilde r_i=\operatorname{LayerNorm}\!\left(r_i+\operatorname{SelfAttn}(R)_i\right),\quad
\alpha_i=\operatorname{softmax}_{i\in\mathrm{real\ news}}
\!\left(\frac{q(m)^\top k(\widetilde r_i)}{\sqrt{64}}\right),
\]
\[
z=\sum_i\alpha_i\widetilde r_i,\qquad
\hat p=\sigma\!\left(\operatorname{MLP}([m;z])\right).
\]
\textbf{Implementation choices.} The residual connection and LayerNorm are our concrete self-attention design choices; they are not established author defaults. Frozen encoders run outside training. Attention weights describe model allocation within a day and do not establish causal market influence.

\sub{Audit completed on 4 October 2026}
\begin{tabularx}{\linewidth}{@{}p{5.1cm}Y@{}}
\toprule
\textbf{Check} & \textbf{Measured outcome}\\\midrule
Source and preparation integrity & 2,111 source returns checked; all 2,107 saved examples match reconstruction; source/processed/cache hashes verified.\\
Feature and model validity & 32,286 finite 384-vectors; valid three-way probabilities; seven integrity tests pass, including padding invariance and attention gradients.\\
Saved experiment replay & 18/18 checkpoints reproduce all 317 test probabilities exactly on CPU; validation losses, selected epochs and saved attention diagnostics agree.\\
Optimization smoke check & Full model fits the first eight training examples: loss 0.6796 to below 0.00001 in 150 steps, with dropout disabled; not evidence of generalization.\\
\bottomrule
\end{tabularx}\par\vspace{3pt}
{\footnotesize Evidence: \code{reports/prototype_audit.json}, \code{reports/results.md}, saved run JSON/checkpoints. Environment: Python 3.12.0, PyTorch 2.5.1+cpu, NumPy 1.26.4. The trained architecture required no repair after this review.}

\newpage
\topic{05 / RESEARCH POSITION AND NEXT STEPS}{What we can defend in the meeting}
\textbf{Claim supported now.} An audited implementation preserves FININ's fusion and attention mechanisms and runs a controlled experiment on substitute data, with baselines and ablations.

\textbf{Claim not supported yet.} Published gains, a forecasting benefit from attention, a trading strategy, and long-memory or market-inefficiency claims remain unverified.

\sub{Research observations that shape the next experiment}
\begin{itemize}
\item \textbf{The paper's primary success criterion is financial performance.} Its accuracy rankings differ from its Sharpe-ratio rankings. Thus a higher raw accuracy on our data would not itself reproduce the central claim [1, Sections 5.5--6.2].
\item \textbf{Lookback length is central to the hypothesis.} FININ aggregates daily market/news representations over up to 20 trading days. Our model tests same-bucket headline interactions but omits the cross-day input history needed to study persistence [1, Eq.\ 7].
\item \textbf{Data substitutions change the available signal.} Generated sentiment, capped coverage, a fixed market description and missing release times differ from TRNA. Current experiments do not isolate their individual effects.
\item \textbf{Retrospective representation is a limitation.} BGE-v1.5 was released in 2023 [3], after our historical test period. Pretraining and sentiment-training provenance have not been audited for historical overlap. This is an offline historical experiment, not a point-in-time simulation.
\end{itemize}

\sub{Recommended next milestone}
\begin{tabularx}{\linewidth}{@{}p{.65cm}>{\raggedright\arraybackslash}p{5.7cm}Y@{}}
\toprule
 & \textbf{Next action} & \textbf{What it would establish}\\\midrule
1 & Freeze the current version; diagnose class collapse using training/validation predictions and learning curves. & Separate optimization or calibration issues from weak usable signal; avoid tuning on the already-inspected test dates.\\
2 & Add 3/5-day inputs, then 10/20 days if feasible; keep a matched mean-pooling baseline and register the comparison before evaluation. & Test whether temporal context changes the value of attention. Use fresh temporal holdouts or a declared rolling-window study.\\
3 & Seek institutional TRNA access or timestamped, instrument-linked news; add a second market. & Move toward the paper's input semantics and test whether findings persist across markets and periods.\\
4 & Specify execution time, costs and risk-free-rate units before adding financial metrics. & Make PnL/Sharpe interpretable and reproducible. Reconcile Appendix B's correctness-based flag and daily rate convention with an executable strategy.\\
\bottomrule
\end{tabularx}\par\vspace{3pt}

\note{\textbf{Discussion point for the professor.} Is the next milestone an architectural proof of concept on accessible data, or a closer empirical replication using licensed news and the original temporal evaluation? The current project is a verified foundation for the first, with a clear path toward the second.}

\sub{References and provenance}
{\footnotesize
[1] M.\ Wang, S.\ B.\ Cohen and T.\ Ma (2024). \emph{Modeling News Interactions and Influence for Financial Market Prediction}. Findings of EMNLP, 3302--3314. \href{https://aclanthology.org/2024.findings-emnlp.189/}{doi:10.18653/v1/2024.findings-emnlp.189}. Reviewed full paper, Figure 2, Eqs.\ 1--9, Tables 1--3 and Appendices A--C.\par
[2] R.\ Saqur. \href{https://huggingface.co/datasets/raeidsaqur/NIFTY}{NIFTY dataset card and files}. Snapshot \code{9b8aef736cba}; separate Yahoo SPY snapshot and checksums recorded in \code{data/spy_manifest.json}.\par
[3] BAAI. \href{https://huggingface.co/BAAI/bge-small-en-v1.5}{BGE-small-en-v1.5 model card}. Snapshot \code{5c38ec7c405e}; 384-dimensional frozen first-token features.\par
[4] mikeysharma. \href{https://huggingface.co/mikeysharma/finance-sentiment-analysis}{Finance sentiment analysis model card}. Snapshot \code{4353ecd593de}; ONNX TinyBERT; positive, neutral and negative outputs reordered and checked with control headlines.\par
Model/dataset cards consulted 4 October 2026. Full revisions and hashes are retained in project manifests. All prototype numbers are derived from local artifacts; no paper score is presented as our measurement.
}
\end{document}
